
\subsection{Group Coursework Component}

\begin{itemize}
  \item When: 2nd half of term 1 (concurrent with classes and lectures).
  \item What: A group project requiring you to develop a piece of software in Python.
  \item Number of submissions: 3
        \begin{itemize}
          \item Your group's codebase and project repository.
          \item Your \href{https://app.ipac-app.education/website/}{IPAC peer-evaluation} review.
          \item Your attendance and performance at the group viva.
        \end{itemize}
\end{itemize}

The coursework component of the module runs concurrently with the lectures and classes, in term 1.
Partway through the year, you will be assigned to a group of 4-5 members, and will be asked to develop a piece of software in Python as detailed in the assignment's document.
You will be expected to develop this software collaboratively with your group, and engage in good development practices as you have seen and learnt from the lectures and course notes.

Your group will submit the final state of your project's repository via moodle.
Markers - both automatic and human - will then examine your repository against a marking rubric to determine how many marks your group will be awarded, based on whether your project meets the specifications set out by the assignment document.
This will give you a score $\codebaseGrade$, awarded to all members of your group, reflecting how well you addressed the functionality requirements of the assignment.
$\codebaseGrade$ will have a maximum value of 50\% (of the overall component's marks).

After the project repository is submitted, each member of your group will also submit an IPAC peer-evaluation, where you will provide feedback to your group members and discuss how well you felt the group worked together.
The IPAC will then calculate an individual weight for you $\ipacWeight$, which is typically in the range $0.8$-$1.2$.
Note that your IPAC weight may be different from that of your group-mates, but your group's weights will always average out to 1.

Finally, your group will be required to attend an in-person viva.
In this viva you will be asked a series of questions about the work you submitted, and your contributions to the software development workflow that your group undertook.
In particular, in your group viva you will be expected to be able to:

\begin{itemize}
  \item Highlight and discuss any contributions you made to design decisions during the project; such as your contribution to issue discussions, or suggestions during your pull-request reviews.
  \item Explain your contributions to the codebase through any pull-requests you authored, as well as evaluate and discuss alternatives that you (or your group) considered.
  \item Justify your software development practices, and any contributions to the repository that reflect this.
\end{itemize}

You will be given an individual score $\vivaGrade$ from your performance in the viva, which will have a maximum value of 50\% (of the overall component's marks).

Your grade for the coursework component $\courseworkGrade$ is then calculated as

\begin{align*}
  \courseworkGrade = (\codebaseGrade \times \ipacWeight) + \vivaGrade,
\end{align*}

being capped at $100\%$ in the event of overflow.
